\section{Exact computation and independently checked certificates}
\label{implementation:certificates}

\subsection{A calculation returns an identity}

The production implementation separates original-graph current
reconstruction, scalar reduction and certificate verification. Its
arithmetic uses exact integers and rational numbers. Finite scalar
tables retain explicit rational coefficients of named constants;
an unresolved scalar period remains an error or an explicitly
unevaluated target. It is never assigned coefficient zero.

For the relative-period algorithm of
Theorem~\ref{periods:finite-span}, fix a cube denominator $U$ with
nonnegative rational coefficients and positive constant coefficient.
Its reciprocal is regular on the entire closed cube. A candidate
target $P/U^q$ is reduced to caller-supplied master forms and an
explicit finite sum of derivatives
\begin{equation}
 \frac{P}{U^q}
 =\sum_j c_j\frac{P_j}{U^{q_j}}
  +\sum_{i,a,r}c_{i,a,r}\,
   \partial_{x_i}\left[x_i(1-x_i)\frac{x^a}{U^r}\right].
 \label{implementation:certificate}
\end{equation}
Here $a$ is a tuple of nonnegative integer powers, $r$ is a
nonnegative integer, and the displayed sums are exactly those
stored in the certificate. All coefficients are rational. The
primitive in each term vanishes on both faces normal to its
differentiation coordinate.

The search enumerates a finite caller-declared range of numerator
degrees and pole orders and performs exact sparse Gaussian
elimination. Each row operation updates the full witness vector.
After a target is eliminated, the implementation clears a common
denominator and checks \eqref{implementation:certificate} as a
polynomial identity. A separate verifier differentiates the stored
primitive expressions with a general symbolic algebra system,
checks each endpoint face, and independently checks that the
denominator, target and masters are the intended ones. This final
binding excludes an algebraically valid certificate for a different
target.

\begin{proposition}[Soundness of a returned target certificate]
Under the stated positivity condition on $U$, a verified identity
\eqref{implementation:certificate} implies the exact corresponding
integral identity on the closed cube.
\end{proposition}
\begin{proof}
Every form and primitive is smooth on the compact cube because
$U$ is bounded below by its positive constant coefficient.
Integrate each derivative in its own coordinate. Its two endpoint
values vanish because of $x_i(1-x_i)$. Fubini's theorem applies
to these continuous functions, so each derivative contributes
zero. Integrating the remaining finite rational identity gives
the result, with no discarded face integral.
\end{proof}

A search failure retains the complete rational residual and its
declared search bounds. It proves neither a nonzero period nor
independence of the supplied masters. Finiteness of the full
relative quotient does not provide a stopping test for an arbitrary
finite relation bank. General effective de Rham algorithms exist
\cite{oaku2001,walther2000}; they are not replaced by a rank plateau
in this implementation.

\subsection{Complete scalar controls}

The three-edge banana graph has two vertices and three parallel
edges, with scalar Kirchhoff polynomial
$U=t_0t_1+t_0t_2+t_1t_2$. At edge powers $(2,2,2)$ and denominator
power three, every one of its six Hepp sectors is
\begin{align}
 \int_0^1\int_0^1\frac{x}{(1+x+xy)^3}\,dy\,dx
 &=\frac12\int_0^1\left(\frac1{(1+x)^2}
                         -\frac1{(1+2x)^2}\right)dx\nonumber\\
 &=\frac12\left(\frac12-\frac13\right)=\frac1{12}.
 \label{implementation:banana}
\end{align}
Thus the complete positive projective period is $6/12=1/2$.
All six ordered sectors have separate complete certificates.
This is a scalar-period control, not an assertion that the
three-edge banana is an admissible original holomorphic graph.

Two further controls test nonrational master values. For the
one-variable denominator $1+x$, the exact antiderivative after
the substitution $u=1+x$ gives
\begin{align}
 \int_0^1\frac{x^7}{(1+x)^5}\,dx
 &=\left[
 \frac{u^3}{3}-\frac{7u^2}{2}+21u-35\log u
 -\frac{35}{u}+\frac{21}{2u^2}-\frac{7}{3u^3}
 +\frac{1}{4u^4}\right]_{u=1}^{u=2}\nonumber\\
 &=\frac{1553}{64}-35\log2.
 \label{implementation:onevar}
\end{align}
Its supplied master forms are $1$ and $(1+x)^{-1}$; the complete
derivative certificate has six nonzero primitives.

For $U=1+xy$, direct integration gives
\begin{align}
 \int_{[0,1]^2}\frac{dx\,dy}{U}&=\frac{\pi^2}{12},\nonumber\\
 \int_{[0,1]^2}\frac{dx\,dy}{U^2}&=\log2,\qquad
 \int_{[0,1]^2}\frac{dx\,dy}{U^3}
 =\frac14+\frac12\log2.\label{implementation:two-base}
\end{align}
For the first identity, the absolutely integrable series of
monomial integrals is
$\sum_{m\ge0}(-1)^m/(m+1)^2=(1-2^{-1})\zeta(2)=\pi^2/12$.
The two remaining identities follow by integrating in $y$ first.
Using $(xy)^2/U^3=U^{-1}-2U^{-2}+U^{-3}$ yields
\begin{equation}
 \int_{[0,1]^2}\frac{x^2y^2}{(1+xy)^3}\,dx\,dy
 =\frac{\pi^2}{12}-\frac32\log2+\frac14.
 \label{implementation:twovar}
\end{equation}
The supplied masters $1$, $1/U$, $x/U$ have values
$1$, $\pi^2/12$, $2\log2-1$; their exact certificate
coefficients are $(-1/2,1,-3/4)$. Appendix~\ref{app:certificates}
prints the complete derivative witnesses. These cube controls
test the reduction engine; their constants are not claimed as
new irreducible periods of the forest family.

\subsection{Scope of the finite census}

The sourced network classification has 1,135 product-forest
topologies through six loops. The finite integration census through
five loops covers 261 topologies at their individually recorded
balanced rational sources, with 5,418 original first-current scalar
entries. Its five-loop part consists of 200 topologies, 200 bare
kernels and 2,200 complete first-current spinors. Those counts are
finite-source results. They do not by themselves interpolate the
arbitrary-source polynomials, which require the separate argument
and calculation of Section~\ref{currents:family}.

The existing scalar tables and original numerator checks are
dependencies of these finite results. The new manuscript replay
distinguishes verification of stored complete certificates from
recomputing every original period table. The former is performed
by the normal replay command; expensive source integrations can
be regenerated by the documented deep replay. No claim that every
period table was independently regenerated during manuscript
preparation is made.

\subsection{Resource bounds and reproducibility}

Graph enumeration, polynomial support, source reconstruction,
relation generation and sparse elimination have explicit finite
budgets. Exceeding a budget raises an error before a partial
result can be returned as a complete target. The scalar tests
exercise divergent weights, exhausted searches, corrupted
witnesses and target substitutions. The complete graph pullbacks
are compared with direct Kirchhoff determinants and their full
Jacobian factors.

The manuscript source and supplement build from the repository
root by running
\begin{center}
\texttt{python3 notes/methods-paper/build.py}
\end{center}
The exact data and manuscript controls replay by
\begin{center}
\texttt{python3 notes/methods-paper/run\_checks.py}
\end{center}
The accepted build records source hashes, page counts, PDF hashes,
compiler version and unresolved-reference/overflow checks. The
verification directory records the independent mathematical and
layout review separately from the numerical certificate results.
